% Options for packages loaded elsewhere
\PassOptionsToPackage{unicode}{hyperref}
\PassOptionsToPackage{hyphens}{url}
\PassOptionsToPackage{dvipsnames,svgnames,x11names}{xcolor}
%
\documentclass[
  11pt]{article}

\usepackage{amsmath,amssymb}
\usepackage{iftex}
\ifPDFTeX
  \usepackage[T1]{fontenc}
  \usepackage[utf8]{inputenc}
  \usepackage{textcomp} % provide euro and other symbols
\else % if luatex or xetex
  \usepackage{unicode-math}
  \defaultfontfeatures{Scale=MatchLowercase}
  \defaultfontfeatures[\rmfamily]{Ligatures=TeX,Scale=1}
\fi
\usepackage{lmodern}
\ifPDFTeX\else  
    % xetex/luatex font selection
\fi
% Use upquote if available, for straight quotes in verbatim environments
\IfFileExists{upquote.sty}{\usepackage{upquote}}{}
\IfFileExists{microtype.sty}{% use microtype if available
  \usepackage[]{microtype}
  \UseMicrotypeSet[protrusion]{basicmath} % disable protrusion for tt fonts
}{}
\makeatletter
\@ifundefined{KOMAClassName}{% if non-KOMA class
  \IfFileExists{parskip.sty}{%
    \usepackage{parskip}
  }{% else
    \setlength{\parindent}{0pt}
    \setlength{\parskip}{6pt plus 2pt minus 1pt}}
}{% if KOMA class
  \KOMAoptions{parskip=half}}
\makeatother
\usepackage{xcolor}
\usepackage[margin=1in]{geometry}
\setlength{\emergencystretch}{3em} % prevent overfull lines
\setcounter{secnumdepth}{-\maxdimen} % remove section numbering
% Make \paragraph and \subparagraph free-standing
\makeatletter
\ifx\paragraph\undefined\else
  \let\oldparagraph\paragraph
  \renewcommand{\paragraph}{
    \@ifstar
      \xxxParagraphStar
      \xxxParagraphNoStar
  }
  \newcommand{\xxxParagraphStar}[1]{\oldparagraph*{#1}\mbox{}}
  \newcommand{\xxxParagraphNoStar}[1]{\oldparagraph{#1}\mbox{}}
\fi
\ifx\subparagraph\undefined\else
  \let\oldsubparagraph\subparagraph
  \renewcommand{\subparagraph}{
    \@ifstar
      \xxxSubParagraphStar
      \xxxSubParagraphNoStar
  }
  \newcommand{\xxxSubParagraphStar}[1]{\oldsubparagraph*{#1}\mbox{}}
  \newcommand{\xxxSubParagraphNoStar}[1]{\oldsubparagraph{#1}\mbox{}}
\fi
\makeatother

\usepackage{color}
\usepackage{fancyvrb}
\newcommand{\VerbBar}{|}
\newcommand{\VERB}{\Verb[commandchars=\\\{\}]}
\DefineVerbatimEnvironment{Highlighting}{Verbatim}{commandchars=\\\{\}}
% Add ',fontsize=\small' for more characters per line
\usepackage{framed}
\definecolor{shadecolor}{RGB}{241,243,245}
\newenvironment{Shaded}{\begin{snugshade}}{\end{snugshade}}
\newcommand{\AlertTok}[1]{\textcolor[rgb]{0.68,0.00,0.00}{#1}}
\newcommand{\AnnotationTok}[1]{\textcolor[rgb]{0.37,0.37,0.37}{#1}}
\newcommand{\AttributeTok}[1]{\textcolor[rgb]{0.40,0.45,0.13}{#1}}
\newcommand{\BaseNTok}[1]{\textcolor[rgb]{0.68,0.00,0.00}{#1}}
\newcommand{\BuiltInTok}[1]{\textcolor[rgb]{0.00,0.23,0.31}{#1}}
\newcommand{\CharTok}[1]{\textcolor[rgb]{0.13,0.47,0.30}{#1}}
\newcommand{\CommentTok}[1]{\textcolor[rgb]{0.37,0.37,0.37}{#1}}
\newcommand{\CommentVarTok}[1]{\textcolor[rgb]{0.37,0.37,0.37}{\textit{#1}}}
\newcommand{\ConstantTok}[1]{\textcolor[rgb]{0.56,0.35,0.01}{#1}}
\newcommand{\ControlFlowTok}[1]{\textcolor[rgb]{0.00,0.23,0.31}{\textbf{#1}}}
\newcommand{\DataTypeTok}[1]{\textcolor[rgb]{0.68,0.00,0.00}{#1}}
\newcommand{\DecValTok}[1]{\textcolor[rgb]{0.68,0.00,0.00}{#1}}
\newcommand{\DocumentationTok}[1]{\textcolor[rgb]{0.37,0.37,0.37}{\textit{#1}}}
\newcommand{\ErrorTok}[1]{\textcolor[rgb]{0.68,0.00,0.00}{#1}}
\newcommand{\ExtensionTok}[1]{\textcolor[rgb]{0.00,0.23,0.31}{#1}}
\newcommand{\FloatTok}[1]{\textcolor[rgb]{0.68,0.00,0.00}{#1}}
\newcommand{\FunctionTok}[1]{\textcolor[rgb]{0.28,0.35,0.67}{#1}}
\newcommand{\ImportTok}[1]{\textcolor[rgb]{0.00,0.46,0.62}{#1}}
\newcommand{\InformationTok}[1]{\textcolor[rgb]{0.37,0.37,0.37}{#1}}
\newcommand{\KeywordTok}[1]{\textcolor[rgb]{0.00,0.23,0.31}{\textbf{#1}}}
\newcommand{\NormalTok}[1]{\textcolor[rgb]{0.00,0.23,0.31}{#1}}
\newcommand{\OperatorTok}[1]{\textcolor[rgb]{0.37,0.37,0.37}{#1}}
\newcommand{\OtherTok}[1]{\textcolor[rgb]{0.00,0.23,0.31}{#1}}
\newcommand{\PreprocessorTok}[1]{\textcolor[rgb]{0.68,0.00,0.00}{#1}}
\newcommand{\RegionMarkerTok}[1]{\textcolor[rgb]{0.00,0.23,0.31}{#1}}
\newcommand{\SpecialCharTok}[1]{\textcolor[rgb]{0.37,0.37,0.37}{#1}}
\newcommand{\SpecialStringTok}[1]{\textcolor[rgb]{0.13,0.47,0.30}{#1}}
\newcommand{\StringTok}[1]{\textcolor[rgb]{0.13,0.47,0.30}{#1}}
\newcommand{\VariableTok}[1]{\textcolor[rgb]{0.07,0.07,0.07}{#1}}
\newcommand{\VerbatimStringTok}[1]{\textcolor[rgb]{0.13,0.47,0.30}{#1}}
\newcommand{\WarningTok}[1]{\textcolor[rgb]{0.37,0.37,0.37}{\textit{#1}}}

\providecommand{\tightlist}{%
  \setlength{\itemsep}{0pt}\setlength{\parskip}{0pt}}\usepackage{longtable,booktabs,array}
\usepackage{calc} % for calculating minipage widths
% Correct order of tables after \paragraph or \subparagraph
\usepackage{etoolbox}
\makeatletter
\patchcmd\longtable{\par}{\if@noskipsec\mbox{}\fi\par}{}{}
\makeatother
% Allow footnotes in longtable head/foot
\IfFileExists{footnotehyper.sty}{\usepackage{footnotehyper}}{\usepackage{footnote}}
\makesavenoteenv{longtable}
\usepackage{graphicx}
\makeatletter
\newsavebox\pandoc@box
\newcommand*\pandocbounded[1]{% scales image to fit in text height/width
  \sbox\pandoc@box{#1}%
  \Gscale@div\@tempa{\textheight}{\dimexpr\ht\pandoc@box+\dp\pandoc@box\relax}%
  \Gscale@div\@tempb{\linewidth}{\wd\pandoc@box}%
  \ifdim\@tempb\p@<\@tempa\p@\let\@tempa\@tempb\fi% select the smaller of both
  \ifdim\@tempa\p@<\p@\scalebox{\@tempa}{\usebox\pandoc@box}%
  \else\usebox{\pandoc@box}%
  \fi%
}
% Set default figure placement to htbp
\def\fps@figure{htbp}
\makeatother

\usepackage{fancyhdr}
\usepackage{setspace}
\onehalfspacing
\pagestyle{fancy}
\lhead{DATA 621}
\rhead{Naomi Buell, Richie Rivera, Alexander Simon, Said Naqwe, and Nana Frimpong}
\setcounter{secnumdepth}{7}
\makeatletter
\@ifpackageloaded{caption}{}{\usepackage{caption}}
\AtBeginDocument{%
\ifdefined\contentsname
  \renewcommand*\contentsname{Table of contents}
\else
  \newcommand\contentsname{Table of contents}
\fi
\ifdefined\listfigurename
  \renewcommand*\listfigurename{List of Figures}
\else
  \newcommand\listfigurename{List of Figures}
\fi
\ifdefined\listtablename
  \renewcommand*\listtablename{List of Tables}
\else
  \newcommand\listtablename{List of Tables}
\fi
\ifdefined\figurename
  \renewcommand*\figurename{Figure}
\else
  \newcommand\figurename{Figure}
\fi
\ifdefined\tablename
  \renewcommand*\tablename{Table}
\else
  \newcommand\tablename{Table}
\fi
}
\@ifpackageloaded{float}{}{\usepackage{float}}
\floatstyle{ruled}
\@ifundefined{c@chapter}{\newfloat{codelisting}{h}{lop}}{\newfloat{codelisting}{h}{lop}[chapter]}
\floatname{codelisting}{Listing}
\newcommand*\listoflistings{\listof{codelisting}{List of Listings}}
\makeatother
\makeatletter
\makeatother
\makeatletter
\@ifpackageloaded{caption}{}{\usepackage{caption}}
\@ifpackageloaded{subcaption}{}{\usepackage{subcaption}}
\makeatother

\usepackage{bookmark}

\IfFileExists{xurl.sty}{\usepackage{xurl}}{} % add URL line breaks if available
\urlstyle{same} % disable monospaced font for URLs
\hypersetup{
  pdftitle={Classification Metrics},
  pdfauthor={Naomi Buell; Richie Rivera; Alexander Simon; Nana Frimpong; Said Naqwe},
  colorlinks=true,
  linkcolor={blue},
  filecolor={Maroon},
  citecolor={Blue},
  urlcolor={Blue},
  pdfcreator={LaTeX via pandoc}}


\title{Classification Metrics}
\usepackage{etoolbox}
\makeatletter
\providecommand{\subtitle}[1]{% add subtitle to \maketitle
  \apptocmd{\@title}{\par {\large #1 \par}}{}{}
}
\makeatother
\subtitle{Homework \#2}
\author{Naomi Buell \and Richie Rivera \and Alexander Simon \and Nana
Frimpong \and Said Naqwe}
\date{2026-03-22}

\begin{document}
\maketitle


\begin{Shaded}
\begin{Highlighting}[]
\FunctionTok{library}\NormalTok{(tidyverse)}
\end{Highlighting}
\end{Shaded}

\begin{verbatim}
-- Attaching core tidyverse packages ------------------------ tidyverse 2.0.0 --
v dplyr     1.1.4     v readr     2.1.5
v forcats   1.0.0     v stringr   1.5.1
v ggplot2   3.5.0     v tibble    3.2.1
v lubridate 1.9.3     v tidyr     1.3.1
v purrr     1.0.2     
-- Conflicts ------------------------------------------ tidyverse_conflicts() --
x dplyr::filter() masks stats::filter()
x dplyr::lag()    masks stats::lag()
i Use the conflicted package (<http://conflicted.r-lib.org/>) to force all conflicts to become errors
\end{verbatim}

\begin{Shaded}
\begin{Highlighting}[]
\FunctionTok{library}\NormalTok{(skimr)}
\FunctionTok{library}\NormalTok{(corrplot)}
\end{Highlighting}
\end{Shaded}

\begin{verbatim}
corrplot 0.95 loaded
\end{verbatim}

\begin{Shaded}
\begin{Highlighting}[]
\FunctionTok{library}\NormalTok{(ggplot2)}
\FunctionTok{library}\NormalTok{(GGally)}
\end{Highlighting}
\end{Shaded}

\begin{verbatim}
Registered S3 method overwritten by 'GGally':
  method from   
  +.gg   ggplot2
\end{verbatim}

\begin{Shaded}
\begin{Highlighting}[]
\FunctionTok{library}\NormalTok{(janitor)}
\end{Highlighting}
\end{Shaded}

\begin{verbatim}

Attaching package: 'janitor'

The following objects are masked from 'package:stats':

    chisq.test, fisher.test
\end{verbatim}

\section{Homework Solutions}\label{homework-solutions}

\begin{enumerate}
\def\labelenumi{\arabic{enumi}.}
\item
  \emph{Download the classification output data set (attached in
  Blackboard to the assignment).}

\begin{Shaded}
\begin{Highlighting}[]
\NormalTok{classifications\_output\_data }\OtherTok{\textless{}{-}} \FunctionTok{read.csv}\NormalTok{(}\StringTok{"classification{-}output{-}data (2).csv"}\NormalTok{)}

\NormalTok{classifications\_output\_data }\SpecialCharTok{|\textgreater{}} \FunctionTok{skim}\NormalTok{()}
\end{Highlighting}
\end{Shaded}

  \begin{longtable}[]{@{}ll@{}}
  \caption{Data summary}\tabularnewline
  \toprule\noalign{}
  \endfirsthead
  \endhead
  \bottomrule\noalign{}
  \endlastfoot
  Name & classifications\_output\_da\ldots{} \\
  Number of rows & 181 \\
  Number of columns & 11 \\
  \_\_\_\_\_\_\_\_\_\_\_\_\_\_\_\_\_\_\_\_\_\_\_ & \\
  Column type frequency: & \\
  numeric & 11 \\
  \_\_\_\_\_\_\_\_\_\_\_\_\_\_\_\_\_\_\_\_\_\_\_\_ & \\
  Group variables & None \\
  \end{longtable}

  \textbf{Variable type: numeric}

  \begin{longtable}[]{@{}
    >{\raggedright\arraybackslash}p{(\linewidth - 20\tabcolsep) * \real{0.2000}}
    >{\raggedleft\arraybackslash}p{(\linewidth - 20\tabcolsep) * \real{0.1053}}
    >{\raggedleft\arraybackslash}p{(\linewidth - 20\tabcolsep) * \real{0.1474}}
    >{\raggedleft\arraybackslash}p{(\linewidth - 20\tabcolsep) * \real{0.0737}}
    >{\raggedleft\arraybackslash}p{(\linewidth - 20\tabcolsep) * \real{0.0632}}
    >{\raggedleft\arraybackslash}p{(\linewidth - 20\tabcolsep) * \real{0.0632}}
    >{\raggedleft\arraybackslash}p{(\linewidth - 20\tabcolsep) * \real{0.0632}}
    >{\raggedleft\arraybackslash}p{(\linewidth - 20\tabcolsep) * \real{0.0737}}
    >{\raggedleft\arraybackslash}p{(\linewidth - 20\tabcolsep) * \real{0.0737}}
    >{\raggedleft\arraybackslash}p{(\linewidth - 20\tabcolsep) * \real{0.0737}}
    >{\raggedright\arraybackslash}p{(\linewidth - 20\tabcolsep) * \real{0.0632}}@{}}
  \toprule\noalign{}
  \begin{minipage}[b]{\linewidth}\raggedright
  skim\_variable
  \end{minipage} & \begin{minipage}[b]{\linewidth}\raggedleft
  n\_missing
  \end{minipage} & \begin{minipage}[b]{\linewidth}\raggedleft
  complete\_rate
  \end{minipage} & \begin{minipage}[b]{\linewidth}\raggedleft
  mean
  \end{minipage} & \begin{minipage}[b]{\linewidth}\raggedleft
  sd
  \end{minipage} & \begin{minipage}[b]{\linewidth}\raggedleft
  p0
  \end{minipage} & \begin{minipage}[b]{\linewidth}\raggedleft
  p25
  \end{minipage} & \begin{minipage}[b]{\linewidth}\raggedleft
  p50
  \end{minipage} & \begin{minipage}[b]{\linewidth}\raggedleft
  p75
  \end{minipage} & \begin{minipage}[b]{\linewidth}\raggedleft
  p100
  \end{minipage} & \begin{minipage}[b]{\linewidth}\raggedright
  hist
  \end{minipage} \\
  \midrule\noalign{}
  \endhead
  \bottomrule\noalign{}
  \endlastfoot
  pregnant & 0 & 1 & 3.86 & 3.24 & 0.00 & 1.00 & 3.00 & 6.00 & 15.00 &
  ▇▃▃▁▁ \\
  glucose & 0 & 1 & 118.30 & 30.48 & 57.00 & 99.00 & 112.00 & 136.00 &
  197.00 & ▂▇▅▂▂ \\
  diastolic & 0 & 1 & 71.70 & 11.80 & 38.00 & 64.00 & 70.00 & 78.00 &
  104.00 & ▁▅▇▅▁ \\
  skinfold & 0 & 1 & 19.80 & 15.69 & 0.00 & 0.00 & 22.00 & 32.00 & 54.00
  & ▇▃▆▅▁ \\
  insulin & 0 & 1 & 63.77 & 88.73 & 0.00 & 0.00 & 0.00 & 105.00 & 543.00
  & ▇▂▁▁▁ \\
  bmi & 0 & 1 & 31.58 & 6.66 & 19.40 & 26.30 & 31.60 & 36.00 & 50.00 &
  ▆▇▇▃▁ \\
  pedigree & 0 & 1 & 0.45 & 0.28 & 0.09 & 0.26 & 0.39 & 0.58 & 2.29 &
  ▇▃▁▁▁ \\
  age & 0 & 1 & 33.31 & 11.18 & 21.00 & 24.00 & 30.00 & 41.00 & 67.00 &
  ▇▃▃▁▁ \\
  class & 0 & 1 & 0.31 & 0.47 & 0.00 & 0.00 & 0.00 & 1.00 & 1.00 &
  ▇▁▁▁▃ \\
  scored.class & 0 & 1 & 0.18 & 0.38 & 0.00 & 0.00 & 0.00 & 0.00 & 1.00
  & ▇▁▁▁▂ \\
  scored.probability & 0 & 1 & 0.30 & 0.23 & 0.02 & 0.12 & 0.24 & 0.43 &
  0.95 & ▇▅▂▂▁ \\
  \end{longtable}
\item
  \emph{The data set has three key columns we will use:}
\end{enumerate}

\begin{itemize}
\item
  \emph{\texttt{class}: the actual class for the observation}
\item
  \emph{\texttt{scored.class}: the predicted class for the observation
  (based on a threshold of 0.5)}
\item
  \emph{\texttt{scored.probability}: the predicted probability of
  success for the observation}

  \emph{Use the table() function to get the raw confusion matrix for
  this scored dataset. Make sure you understand the output. In
  particular, do the rows represent the actual or predicted class? The
  columns?}

\begin{Shaded}
\begin{Highlighting}[]
\NormalTok{confusion\_matrix }\OtherTok{\textless{}{-}} \FunctionTok{table}\NormalTok{(}
\NormalTok{classifications\_output\_data}\SpecialCharTok{$}\NormalTok{class,}
\NormalTok{classifications\_output\_data}\SpecialCharTok{$}\NormalTok{scored.class}
\NormalTok{)}
\NormalTok{confusion\_matrix}
\end{Highlighting}
\end{Shaded}

\begin{verbatim}

      0   1
  0 119   5
  1  30  27
\end{verbatim}

\begin{Shaded}
\begin{Highlighting}[]
\NormalTok{tn }\OtherTok{\textless{}{-}}\NormalTok{ confusion\_matrix[}\DecValTok{1}\NormalTok{, }\DecValTok{1}\NormalTok{]}
\NormalTok{fn }\OtherTok{\textless{}{-}}\NormalTok{ confusion\_matrix[}\DecValTok{2}\NormalTok{, }\DecValTok{1}\NormalTok{]}
\end{Highlighting}
\end{Shaded}

  The rows represent the actual class (\texttt{class}) and the columns
  represent the predicted class (\texttt{scored.class}). For example,
  there are 119 observations where the actual class is 0 and the
  predicted class is also 0 (true negatives), and there are 30
  observations where the actual class is 1 but the predicted class is 0
  (false negatives).
\end{itemize}

\begin{enumerate}
\def\labelenumi{\arabic{enumi}.}
\setcounter{enumi}{2}
\item
  \emph{Write a function that takes the data set as a dataframe, with
  actual and predicted classifications identified, and returns the
  accuracy of the predictions.}

  \[
  𝐴𝑐𝑐𝑢𝑟𝑎𝑐𝑦 = \frac{𝑇𝑃 + 𝑇𝑁}{𝑇𝑃 + 𝐹𝑃 + 𝑇𝑁 + 𝐹𝑁}
  \]

\begin{Shaded}
\begin{Highlighting}[]
\NormalTok{calculate\_accuracy }\OtherTok{\textless{}{-}} \ControlFlowTok{function}\NormalTok{(data, actual\_col, predicted\_col) \{}
\NormalTok{confusion\_matrix }\OtherTok{\textless{}{-}} \FunctionTok{table}\NormalTok{(data[[actual\_col]], data[[predicted\_col]])}
\NormalTok{tp }\OtherTok{\textless{}{-}}\NormalTok{ confusion\_matrix[}\DecValTok{2}\NormalTok{, }\DecValTok{2}\NormalTok{]}
\NormalTok{tn }\OtherTok{\textless{}{-}}\NormalTok{ confusion\_matrix[}\DecValTok{1}\NormalTok{, }\DecValTok{1}\NormalTok{]}
\NormalTok{fp }\OtherTok{\textless{}{-}}\NormalTok{ confusion\_matrix[}\DecValTok{1}\NormalTok{, }\DecValTok{2}\NormalTok{]}
\NormalTok{fn }\OtherTok{\textless{}{-}}\NormalTok{ confusion\_matrix[}\DecValTok{2}\NormalTok{, }\DecValTok{1}\NormalTok{]}
\NormalTok{accuracy }\OtherTok{\textless{}{-}}\NormalTok{ (tp }\SpecialCharTok{+}\NormalTok{ tn) }\SpecialCharTok{/}\NormalTok{ (tp }\SpecialCharTok{+}\NormalTok{ fp }\SpecialCharTok{+}\NormalTok{ tn }\SpecialCharTok{+}\NormalTok{ fn)}
\FunctionTok{return}\NormalTok{(accuracy)}
\NormalTok{\}}

\CommentTok{\# Test function}
\FunctionTok{calculate\_accuracy}\NormalTok{(classifications\_output\_data, }\StringTok{"class"}\NormalTok{, }\StringTok{"scored.class"}\NormalTok{)}
\end{Highlighting}
\end{Shaded}

\begin{verbatim}
[1] 0.8066298
\end{verbatim}
\item
  \emph{Write a function that takes the data set as a dataframe, with
  actual and predicted classifications identified, and returns the
  classification error rate of the predictions.}

  \[
  𝐶𝑙𝑎𝑠𝑠𝑖𝑓𝑖𝑐𝑎𝑡𝑖𝑜𝑛 𝐸𝑟𝑟𝑜𝑟 𝑅𝑎𝑡𝑒 = \frac{𝐹𝑃 + 𝐹𝑁}{𝑇𝑃 + 𝐹𝑃 + 𝑇𝑁 + 𝐹𝑁}
  \]

  \emph{Verify that you get an accuracy and an error rate that sums to
  one.}
\item
  \emph{Write a function that takes the data set as a dataframe, with
  actual and predicted classifications identified, and returns the
  precision of the predictions.}

  \[
  𝑃𝑟𝑒𝑐𝑖𝑠𝑖𝑜𝑛 = \frac{𝑇𝑃}{𝑇𝑃 + 𝐹𝑃}
  \]
\item
  \emph{Write a function that takes the data set as a dataframe, with
  actual and predicted classifications identified, and returns the
  sensitivity of the predictions. Sensitivity is also known as recall.}

  \[
  𝑆𝑒𝑛𝑠𝑖𝑡𝑖𝑣𝑖𝑡𝑦 = \frac{𝑇𝑃}{𝑇𝑃 + 𝐹𝑁}
  \]
\item
  \emph{Write a function that takes the data set as a dataframe, with
  actual and predicted classifications identified, and returns the
  specificity of the predictions.}

  \[
  𝑆𝑝𝑒𝑐𝑖𝑓𝑖𝑐𝑖𝑡𝑦 = \frac{𝑇𝑁}{𝑇𝑁 + 𝐹𝑃}
  \]
\item
  \emph{Write a function that takes the data set as a dataframe, with
  actual and predicted classifications identified, and returns the F1
  score of the predictions.}

  \[
  𝐹1 𝑆𝑐𝑜𝑟𝑒 = \frac{2 × 𝑃𝑟𝑒𝑐𝑖𝑠𝑖𝑜𝑛 × 𝑆𝑒𝑛𝑠𝑖𝑡𝑖𝑣𝑖𝑡𝑦}{𝑃𝑟𝑒𝑐𝑖𝑠𝑖𝑜𝑛 + 𝑆𝑒𝑛𝑠𝑖𝑡𝑖𝑣𝑖𝑡𝑦}
  \]
\item
  \emph{Before we move on, let's consider a question that was asked:
  What are the bounds on the F1 score? Show that the F1 score will
  always be between 0 and 1. (Hint: If 0 \textless{} 𝑎 \textless{} 1 and
  0 \textless{} 𝑏 \textless{} 1 then 𝑎𝑏 \textless{} 𝑎.)}
\item
  \emph{Write a function that generates an ROC curve from a data set
  with a true classification column (class in our example) and a
  probability column (scored.probability in our example). Your function
  should return a list that includes the plot of the ROC curve and a
  vector that contains the calculated area under the curve (AUC). Note
  that I recommend using a sequence of thresholds ranging from 0 to 1 at
  0.01 intervals.}
\item
  \emph{Use your created R functions and the provided classification
  output data set to produce all of the classification metrics discussed
  above.}
\item
  \emph{Investigate the caret package. In particular, consider the
  functions confusionMatrix, sensitivity, and specificity. Apply the
  functions to the data set. How do the results compare with your own
  functions?}
\item
  \emph{Investigate the pROC package. Use it to generate an ROC curve
  for the data set. How do the results compare with your own functions?}
\end{enumerate}

\section{Classification Metrics}\label{classification-metrics}

\emph{Upon following the instructions below, use your created R
functions and the other packages to generate the classification metrics
for the provided data set.}




\end{document}
